\chapter{Conclusion and Future Work}
\label{ch:conclusion}

\section{Summary of Contributions}

This dissertation has presented AIEF as an exploratory instrument for a contested ethics-rule space: free-text proposals in, structured and authority-formatted assessments out, with the rules and priorities that produce those assessments available for inspection and revision. The system integrates an OWL ontology harmonising 207 requirements across REAMS, the EU AI Act, Horizon Europe, and ACM/NeurIPS; a knowledge graph encoding these requirements alongside 70 real-world AI incidents with 342 Charter rights mappings; a retrieval-augmented generation pipeline; a multi-authority documentation generator; and a SHACL risk-dimension layer with harm-evidenced, evaluator-editable priorities.

The evaluation answers the three findings-oriented research questions. For RQ1, landscape analysis shows eight Charter rights covered by all four frameworks, Article~24 covered by REAMS alone, and four rights with incident evidence but zero requirement mappings after dual-IRI cleanup (Articles~13, 15, 17, and~48). For RQ2, hybrid RAG and SHACL trade free-text coverage and narrative explanation against determinism and editability; the ablation study confirms that LLM-alone baselines lack requirement-level provenance. For RQ3, one assessment yields mean documentation completeness of 100\% (FRIA, ACM/NeurIPS), 86\% (REAMS), and 61\% (Horizon Europe), with mean fundamental-rights coverage of 0.458 against reference annotations --- a stricter primary metric than the secondary 90\% risk-label accuracy, by design; risk accuracy is retained as secondary. Retrieval recall remains 0.99 across model configurations.

AIEF realises the automated compliance assessment tool that Pandit and Rintam\"{a}ki \cite{pandit2024fria} explicitly proposed as future work in their FRIA ontology research. To the best of our knowledge, no existing tool combines free-text proposal input, multi-framework coverage, an OWL ontology with incident grounding, EU Charter rights mapping, risk classification, and requirement-level traceability within a single RAG-based architecture; AIEF's contribution is the integration of these components into an operational, evaluated system rather than any one component in isolation.


\section{Practical Implications}

The system has three immediate practical applications. For \textbf{researchers}, AIEF can serve as a pre-submission screening tool, identifying ethics requirements and potential risks before formal review, reducing the likelihood of incomplete applications and iterative resubmissions. For \textbf{ethics committees}, the system can provide a preliminary assessment that supplements (but does not replace) human review, highlighting areas requiring closer attention. For \textbf{institutions}, the system's multi-framework coverage enables consistent application of ethics requirements across research groups, reducing the variability that arises from individual researchers' differing awareness of regulatory obligations.


\section{Future Work}

Several directions for future research emerge from this work. Two items originally listed as proposed are now demonstrated: scrutable SHACL risk prioritisation with harm-evidenced defaults (\S\ref{sec:shacl}), and multi-authority documentation generation from a single assessment (\S\ref{sec:docgen}). Remaining directions:

\textbf{Framework extension.} The ontology architecture supports the addition of further governance frameworks. ISO/IEC 42001 (AI Management System), the NIST AI Risk Management Framework, and sector-specific regulations (e.g., FDA guidance for clinical AI, financial services AI regulations) could be integrated by defining new requirement individuals and Charter mappings within the existing schema.

\textbf{Knowledge graph maintenance and gap filling.} The AIAAIC Repository and AI Incident Database are continuously updated. An automated or semi-automated pipeline for ingesting new incidents, annotating them with Charter rights mappings, and inserting them into the knowledge graph would ensure the system's evidence base remains current. Targeted incident expansion for the four harm-without-governance rights identified in \S\ref{sec:landscape} would strengthen those gap claims.

\textbf{Embedding-based retrieval.} The current keyword-based SPARQL retrieval could be augmented with semantic embedding similarity to improve recall on proposals that use terminology not directly matching requirement text --- and to lift fundamental-rights coverage (\S\ref{sec:rights-coverage}) above the current mean of 0.458.

\textbf{Expert validation.} Completing the planned expert review with ADAPT Centre researchers would provide inter-rater reliability metrics (Spearman $\rho$) for system output quality, following the PASTA \cite{pasta2026} evaluation methodology.

\textbf{Namespace separation.} The current unified namespace (\texttt{https://w3id.org/aief/}) could be separated into per-framework namespaces to support modular ontology publication and independent versioning, as recommended during the expert ontology review. Dual-IRI aliases for Articles~6 and~7 were merged before the landscape figures in \S\ref{sec:landscape}; remaining naming hygiene across supporting exports should accompany that work.

\textbf{OOPS! remediation.} The OOPS! pitfall scan (\S\ref{sec:oops-validation}) identified concrete, addressable gaps: domain/range declarations for five data properties, \texttt{rdfs:label}/\texttt{rdfs:comment} annotations for 78 supporting-vocabulary elements, and \texttt{owl:inverseOf} declarations for 17 object properties. None require redesigning the class hierarchy.

\textbf{Retrieval precision.} SPARQL retrieval currently prioritises recall over precision by design, retrieving a broad candidate set (typically 150--200 of 207 requirements per proposal) for the LLM to filter; this keeps context recall high (mean 0.80 across the 20 synthetic proposals per the RAGAS-style evaluation in \S\ref{sec:ragas-validation}) but yields low context precision. Narrowing retrieval — for example through embedding similarity re-ranking before LLM hand-off — is a natural next step now that this trade-off has been measured directly rather than assumed.

\textbf{Stakeholder participation.} Following \cite{jarvers2026engaged,ullstein2026pafria}, investigating how AIEF's automated screen could feed into, rather than substitute for, a participatory review process involving affected stakeholders and internal experts.

\textbf{Corporate and regulatory extension.} While this dissertation focuses on research proposal assessment, the architecture is designed to support extension to corporate AI governance contexts. Replacing REAMS requirements with corporate AI governance policies and adding sector-specific regulatory requirements would enable the system to serve as a compliance screening tool for industry AI deployments.

\textbf{Fine-tuning.} The current system uses general-purpose LLMs with prompt engineering. Fine-tuning a model on ethics assessment outputs could improve generation F1 scores, particularly for smaller models, reducing the dependency on large parameter counts.

\textbf{SHACL rule-base expansion.} The six-shape demonstrator (\S\ref{sec:shacl}) establishes the architecture; populating a fuller rule base --- and deciding whose priorities the defaults encode --- remains open community work rather than unfinished dissertation scope.


\section{Concluding Remarks}

The gap between the growing body of AI ethics governance requirements and the tools available to researchers for navigating them represents both a practical problem and a research opportunity. This dissertation has demonstrated that an ontology-driven knowledge graph, combined with retrieval-augmented generation, multi-authority documentation projection, and scrutable SHACL prioritisation, can make that contested rule space explorable rather than opaque. The landscape gaps, architecture trade-offs, and documentation completeness figures are the findings; the 0.99 retrieval recall confirms that the knowledge engineering underneath them is robust and reusable. As both the regulatory landscape and LLM capabilities continue to evolve, the modular architecture of AIEF provides a foundation that can adapt to both.
